\uuid{t0YA}
\titre{Image directe}
\auteur{bodin ; hanani ; mzari}
\organisation{exo7}
\licence{CC BY-NC-SA 4.0}
\datecreate{2020-01-20}
\niveau{L1}
\module{Algèbre}
\chapitre{Logique, ensemble, raisonnement}
\sousChapitre{Ensemble}
\difficulte{2}
\theme{Ensembles, applications, image directe}
\exo7qcmid{lille-1.ensembles.016}

\contenu{
\begin{question}
\qclassification{100.02,101.01,102.02,102.02}
\qtags{motcle=image directe}
On considère l'application $f:\Rr\to \Rr$ définie par
$$\forall x\in \Rr,\; f(x)=x^2+1.$$
Quelle est la bonne réponse ?
\begin{answers}
    \bad{$f(\Rr)=\Rr$}
    \bad{$f(\Rr)=[0,+\infty [$}
    \bad{$f(\Rr)=]1,+\infty [$}
    \good{$f(\Rr)=[1,+\infty [$}
\end{answers}
\begin{explanations}
Pour tout $x\in \Rr$, $f(x)\geq 1$. Donc $f(\Rr)\subset [1,+\infty[$. Réciproquement, tout $y\in [1,\infty [$ admet un antécédent. Donc $[1,+\infty[ \subset f(\Rr)$.
\end{explanations}
\end{question}
}
